% main-causal.tex — Dempsey Lab manuscript template (JASA-style article)
% ESTIMAND-DRIVEN CAUSAL structure: Setup and Estimand / Identification /
% Estimation. For general methods papers use main.tex instead (Intro w/
% Contributions, Related Work, Method: Notation / Core method / Theory).
% Keep one, delete the other, and name the keeper main.tex. Comments
% explain what belongs in each section.
\documentclass[12pt]{article}

\usepackage{amsmath, amssymb, mathtools}
\usepackage{amsthm}
\usepackage[toc,page]{appendix}
\usepackage{graphicx}
\usepackage{placeins}
\usepackage{pgfplots}
\usepackage[shortlabels]{enumitem}
\usepackage{xcolor}
\usepackage{natbib}
\usepackage{float}
\usepackage[parfill]{parskip}
\usepackage{url}
\usepackage{hyperref}
\usepackage{booktabs}
\usepackage{bbm, bm}
\usepackage{amsfonts}

% Shared lab notation (operators, sequential-decision macros, theorem
% environments). Load BEFORE the newtx fonts. Do not redefine its macros.
\usepackage{lab-notation}

\usepackage{newtxtext}
\usepackage[subscriptcorrection]{newtxmath}

% DON'T change margins - should be 1 inch all around.
\addtolength{\oddsidemargin}{-.5in}%
\addtolength{\evensidemargin}{-.5in}%
\addtolength{\textwidth}{1in}%
\addtolength{\textheight}{-.3in}%
\addtolength{\topmargin}{-.8in}%

% NOTE: To produce blinded version, replace "0" with "1" below.
\newcommand{\blind}{0}

\def\spacingset#1{\renewcommand{\baselinestretch}{#1}\small\normalsize}
\spacingset{1}

\pgfplotsset{compat=1.18}

\begin{document}

\if1\blind
{
  \title{\bf [Project Title]}
  \author{Author 1\thanks{The authors gratefully acknowledge relevant funding sources in the unblinded version.}\\
    Department of Statistics}
  \maketitle
} \fi

\if0\blind
{
  \bigskip
  \bigskip
  \bigskip
  \begin{center}
    {\LARGE\bf [Project Title]}

    \medskip
    S. Ferlito and W. Dempsey

    Department of Biostatistics, University of Michigan, Ann Arbor, U.S.A.
  \end{center}
  \medskip
} \fi

\begin{abstract}
% One paragraph, written to be readable by someone who reads nothing else:
% the setting, the gap, the proposed estimand/method, what theory is
% established, and one sentence each on simulations and the application.
[Abstract placeholder.]
\end{abstract}

\noindent%
{\it Keywords:} [keyword 1]; [keyword 2]; [keyword 3].
\vfill

\newpage
\spacingset{1.45} % DON'T change the spacing!

\section{Introduction}
\label{sec:intro}
% What the scientific setting is, what question existing methods answer,
% and precisely why they fail for this one. End with a numbered list of
% contributions and a roadmap of the paper. Cite generously — this section
% positions the work relative to adjacent literatures.
[Placeholder.]

\section{Setup and Estimand}
\label{sec:setup}
% Define the observed data structure, potential outcomes, and the estimand
% BEFORE any method details. State causal assumptions as numbered
% \begin{assumption} environments. The reader should finish this section
% knowing exactly what quantity the paper targets and under what conditions
% it is well-defined.
Let $X_0$ denote a vector of baseline covariates. At the beginning of Stage 1, each participant is randomly assigned to a first-stage intervention $Z_1 \in \{-1, 1\}$. The known Stage 1 randomization probabilities are $P(Z_1 = z_1 | X_0)$. During Stage 1, participants are micro-randomized at $t^* -1$ decision points. That is, participants are randomized to an intervention to an intervention $A_t \in \{0,1\}$ at each time point $t = 1, \cdots, t^* - 1$. All participants transition to Stage 2 at time $t = t^*$ and are randomized to a second-stage intervention $Z_2$. The possible values for $Z_2$ and the known Stage 2 randomization probabilities depend on the embedded SMART design. 

\quad Unrestricted SMART: $Z_2 \in \{0,1\}$, $P(Z_2 = z_2)$

\quad Restricted SMART, non-responders re-randomized: $Z_2 \in \{-1, 0, 1\}$, $P(Z_2 = z_2 | R)$

\quad Restricted SMART, non-responders to $Z_1 = 1$ re-randomized: $Z_2 \in \{-1, 0, 1\}$, $P(Z_2 = z_2 | R, Z_1)$ \textcolor{blue}{(technically should condition on $Z_1$)}

Participants are micro-randomized at $T-t^*$ decision points in Stage 2 (such that there are $T$ total decision points for micro-randomization across Stages 1 and 2). That is, participants are randomized to an intervention $A_t \in \{0,1\}$ at each time point $t = t^* + 1, \cdots, T$. Across Stages 1 and 2, we use $X_t$ to denote individual and contextual information collected at the $t$-th time point, which is measured before receiving intervention $A_t$. We let $Y_{t+1}$ denote the proximal response which is observed after receiving intervention $A_t$. The collection of observed data is
\begin{equation*}
    X_0, Z_1, X_1, A_1, Y_2, \cdots Z_T, X_T, A_T, Y_{T+1}
\end{equation*}

We use overbars to denote full history (e.g., $\bar{A} = (A_1, A_2, \ldots, A_T)$). In addition, we let $H_t = (\bar{X}_t, \bar{Z}_t, \bar{A}_t, \bar{Y}_t)$ denote the complete history of observable information up to $t$. The micro-randomization probabilities are also known and are denoted by $P(A_t = a | H_t) = p_t(a_t | H_t)$.

Note that the assigned Stage 1 and Stage 2 interventions are part of the history $H_t$. In AIMS, we compare JITAI plus self-monitoring versus self-monitoring alone. Only participants in the JITAI arm undergo micro-randomization, so that $p_t(a_t | H_t, Z_1 = -1) = 0$.

Note that Sustain is an unrestricted SMART, with three SMART stages (biweekly slow time-scale decision points). In addition, all participants undergo micro-randomization. \textcolor{blue}{If we are using Sustain as a data example but we expect most HEDs to have two SMART stages, should we put two time-point case $Z_1, Z_2$ in main text and three time-point case in supplement?}

Define potential outcomes and DTRs next. 


\begin{assumption}[Consistency, no interference, and sequential randomization]
\label{ass:core}
[Placeholder.]
\end{assumption}

\section{Identification}
\label{sec:id}
% Derive observed-data representations of the estimand under the stated
% assumptions. Each representation is a theorem or proposition; proofs go
% to the appendix. Explain what each representation buys (interpretation,
% estimation strategy, variance properties).
[Placeholder.]

\section{Estimation}
\label{sec:estimation}
% Estimators built from the identification results: the simple/direct
% estimator first, then augmented or cross-fitted refinements. Include the
% large-sample theory (asymptotic normality, variance estimation) here or
% in a dedicated subsection. Practical fitting details (cross-fitting
% scheme, nuisance regressions) belong here, not the appendix.
[Placeholder.]

\section{Simulation Study}
\label{sec:sims}
% One subsection per study, mirroring the simulations/ subdirectories in
% the repo — each answers a single named question (bias? coverage?
% efficiency gain? robustness to X?). State the data-generating process,
% the estimators compared, and the metric before showing tables.
[Placeholder.]

\section{Application}
\label{sec:application}
% The motivating data analysis. Describe the study, why the method fits,
% results, and a sensitivity analysis. Numbers here must be reproducible
% from the code repository.
[Placeholder.]

\section{Discussion}
\label{sec:discussion}
% Limitations, honest scope of the assumptions, and 2-3 concrete
% directions for future work.
[Placeholder.]

\bibliographystyle{agsm}
\bibliography{refs}

\begin{appendices}

\section{Proofs}
\label{app:proofs}
% One appendix section per main-text section that has results, in order.
[Placeholder.]

\section{Additional Simulation and Application Details}
\label{app:extra}
% Extra tables/figures, diagnostics, and anything reviewers may want but
% readers don't need.
[Placeholder.]

\section{Code and Data Availability}
\label{app:code}
% Link the GitHub repository and state where the data live and how access
% is governed (mirror data/README.md).
[Placeholder.]

\end{appendices}

\end{document}
